problem:  
Let \(G\) be a compact, connected Lie group acting smoothly and effectively with cohomogeneity one on a compact, simply connected manifold \(M^n\).  
Assume \(M\) admits a family of \(G\)-invariant Riemannian metrics \(\{g_t\}_{t>0}\) with the following uniform control:  
- the sectional curvature satisfies \(0 \le \sec_{g_t} \le C\) for some constant \(C>0\) independent of \(t\);  
- the total diameter \(\mathrm{diam}(M,g_t)\) is uniformly bounded;  
- for one of the singular orbits \(G/K_-\), the intrinsic diameter \(\mathrm{diam}\big((G/K_-), g_t|_{G/K_-}\big) \to 0\) as \(t\to 0\).  

Under these conditions, the sequence \((M,g_t)\) collapses (in the pointed Gromov–Hausdorff sense) to a limit space \(X\) with \(\dim X < n\).

**Open Problem (Equivariant Cohomogeneity–One Collapse Rigidity):**  
Does the limit space \(X\) necessarily carry a continuous \(G\)-action whose orbit space \(X/G\) is homeomorphic to an interval (i.e., does \(X\) inherit a cohomogeneity-one action)?  
Equivalently, can there exist a collapsing sequence of \(G\)-invariant, nonnegatively curved, cohomogeneity-one manifolds whose Gromov–Hausdorff limit **fails** to support any effective continuous \(G\)-action with one-dimensional quotient?

Why is it a "good" problem:  
This reformulation makes the issue **mathematically precise**—the collapse has bounded curvature, diameter, and a specific orbit degeneration regime—and directly targets the unknown relationship between **symmetry and collapse limits**.  

1. **Conceptually pure yet open:**  
   It asks in the simplest symmetric setting—cohomogeneity one—whether curvature-controlled collapse can destroy the group symmetry in the limit.  No general theorem guarantees the persistence of continuous group actions under collapse with only nonnegative curvature, making this an authentically open question.  

2. **Bridge across fields:**  
   The problem links the theory of **transformation groups** (via \(G\)-actions and orbit structures) with **metric collapse theory** (Cheeger–Fukaya–Gromov, Yamaguchi), and with **Alexandrov geometry** (structure of limits under curvature bounds).  Solving it would likely involve developing new equivariant techniques for Alexandrov convergence, thus uniting multiple major frameworks.  

3. **Simplicity of statement, depth of consequence:**  
   The question’s language—"does the group symmetry survive a controlled collapse?"—is accessible and geometric, yet its resolution may require sophisticated new concepts in equivariant metric geometry, possibly revealing rigidity phenomena or novel counterexamples.  

4. **Forward-looking:**  
   A positive answer would establish a new rigidity law governing the behavior of symmetry under curvature-preserving collapse; a counterexample would expose a mechanism for **symmetry-breaking in the limit**.  Either outcome would significantly advance our understanding of the geometry and topology of nonnegatively curved manifolds with large symmetry.